
\section{ADARA Network Protocol}

\subsection{Overview}

The objective of this protocol is to provide a simple, easy to parse
structure for the transmission and storage of the data associated with
an experimental run at the SNS. It must accomodate high data rate of the
neuton event information from the detector electronics while offering
the flexability to describe and report data from a diverse and dynamic
collection of sample environment sensors and control systems.

While the format is primarily used for network transmission over TCP
streams, certain packets will also be used over UDP connections when
communicating with embedded systems. Other packets are defined to enable
partial recovery from disk errors at inopportune locations in the stream.

Some devices (and/or processes) in the system will only generate or
operate on a subset of the defined packet formats. For example, it is
expected that the Data Stream Packetizers will generate Raw Event packets,
which will be later transformed into Neutron Event packets once the
physical to logical mapping is performed. To this end, the protocol is
designed allow parsers to skip over packets that are unknown or irrelevant
to the task at hand.

\subsection{Common Protocol Information}

As the bulk of computing power at the SNS is based on the Intel architecture,
all multi-byte values in the protocol will be represeneted in little-endian
format to avoid the overhead of byte-swaping every value. Any big-endian
hardware added in the future will need to accomodate the existing format.

All packets must be a multiple of 4 bytes in length to keep the integer
fields naturally aligned.

The common field types are given in Table~\ref{table:protocol_field_types}.

\begin{table}[h]
  \begin{center}
    \begin{tabular}{l | l}
	Type & Definition \\
	\hline
	{\bf byte} & 8 bit value, unspecified sign \\
	{\bf word} & 32 bit value, unspecified sign \\
	{\bf u16} & 16 bit unsigned integer \\
	{\bf s32} & 32 bit signed integer \\
	{\bf u32} & 32 bit unsigned integer \\
	{\bf u64} & 64 bit unsigned integer \\
	{\bf u128} & 128 bit unsigned integer \\
	{\bf double} & IEEE 754 64-bit floating point number \\
	{\bf string} & UTF-8 encoded character data of specified length \\
    \end{tabular}
  \end{center}
  \caption {Field Type Definitions}
  \label{table:protocol_field_types}
\end{table}

\subsection{Common Packet Header}

\begin{figure}
  \centering
  \begin{bytefield}{32}
    \bitheader{31,0} \\
    \begin{leftwordgroup}{Header}
      \wordbox{1}{Payload Length (u32)} \\
      \wordbox{1}{Format (u32)} \\
      \wordbox{1}{Timestamp (seconds, u32)} \\
      \wordbox{1}{Timestamp (nanoseconds, u32)}
    \end{leftwordgroup} \\
    \begin{leftwordgroup}{Payload}
      \wordbox[lrt]{1}{Data} \\
      \skippedwords \\
      \wordbox[lrb]{1}{}
    \end{leftwordgroup}
  \end{bytefield}
  \caption{ADARA Network Packet Structure}
  \label{fig:protocol_packet}
\end{figure}

\begin{figure}
  \centering
  \begin{bytefield}{32}
    \bitheader{31,8,7,0} \\
    \bitbox{24}{Format ID} &
    \bitbox{8}{Version}
  \end{bytefield}
  \caption{Packet Format Fields}
  \label{fig:protocol_format_fields}
\end{figure}


To ensure that parsers are able to find the next packet in the stream, each
packet includes a 16 byte header. This header consists of a payload length,
packet type, and timestamp information as seen in
Figure~\ref{fig:protocol_packet}.

Each packet has an option payload. If the packet carries information beyond the
header itself, then the {\bf Payload Length} field will be non-zero and a
mulitple of 4. It is expected that stream parsers will first read the full 16
bytes of the header. The {\bf Payload Length} field will then give the length
remaining to be read from the stream for the packet.

Following the {\bf Payload Length}, there is a {\bf Format} field. This 4 byte
field allows for over 4 billion packet types when used sequentially, though the
field is broken into subfields. While the {\bf Format} field may be used
as a flat integer by the parser, the subfields allow some developer
convienience by keeping different version of packets with the same semantics
close together numberically. The field breakout is shown in
Figure~\ref{fig:protocol_format_fields}.

Once a {\bf Format} value has been deployed to production,
the packet format associated with that number shall become immutable. Should
changes be needed to the physical layout or semantic content of that packet, a
new packet type must be allocated and defined in this document. This will
ensure that tools are able to recognize when they are operating on data that
they do not understand.

\todo{Build table of packet types}{%
Once a decision is reached regarding the allocation policy of
packet types, we must assign values to the packet formats described
in this document, and should include them in a table for easy reference.
Bonus points if we have a \LaTeX command to automatically populate the
table and check for duplicates. There should be distinct ranges allocated
for detector system vs ADARA packets, though this is for convienience and
should not carry semantic meaning. For now, I suggest 0-8191 for detector
and 16384+ for ADARA.}

Following the {\bf Format} field, there is a 64-bit {\bf Timestamp}, broken
into two 32-bit fields. The timestamp gives the seconds (first field) and
nanoseconds since the EPICS epoch, defined as midnight, January 1, 1990.  All
timestamp values are given in UTC time.  This field shall be filled with the
pulse ID given by the accelerator system for all data packets carrying neutron
or ``fast'' environment data. As this field will be used to correlate data from
different sources with the neutron event data, all other packets shall use as
accurate a clock as is available on the local system. This clock should be
closely synchronized to same reference source as the accelerator system.

\subsection{Detector System generated packets}

The following packets are generated by the detector system. As such, they will
use the Pulse ID from the Real Time Data Link (RTDL) system to identify each
event an accelerator pulse. Note that while the time-of-flight information for
each pixel has been corrected for frame skew~\footnote{Frame skew indicates a
time-of-flight larger than the intra-pulse time.  For more information on frame
skew, please see the document ``Neutron Event Timing''.}, the ``fast'' metadata
information has no correction applied.

\todo{Define Detector Control Packets}{%
We have currently defined the packets required to communicate event data from
the detectors to the rest of the aquisition system, but have not defined the
packets that will be used to initialize and control the detector system, as
they are not believed to block progress and the software developers do not
currently have the knowledge required to define them. The detector system
developers will fill them in at a later date.}

\newpage
\subsubsection{Detector System: Raw event data}
\label{section:protocol_raw_event_packet}

\begin{figure}[h]
  \centering
  \begin{bytefield}{32}
    \bitheader{31,0} \\
    \wordbox{1}{Length $8 N + 16$} \\
    \wordbox{1}{TBD} \\
    \wordbox{1}{Pulse ID (seconds)} \\
    \wordbox{1}{Pulse ID (nanoseconds)} \\

    \bitheader{31,30,24,23,14,13,0} \\
    \wordbox{1}{Source ID} \\
    \bitbox{1}{\tiny E \\ O \\ P} &
    \bitbox{7}{Reserved} &
    \bitbox{10}{Cycle} &
    \bitbox{14}{Pkt Seq} \\
    \wordbox{1}{Pulse Charge (units of 10 pC, u32)} \\
    \wordbox{1}{Pulse Energy (eV?, u32)} \\
    \wordbox{2}{Event 1} \\
    \wordbox{2}{$\cdots$} \\
    \wordbox{2}{Event $N$} \\
  \end{bytefield}
  \caption{Packet Type TBD: Raw Event Data}
  \label{fig:protocol_packet_raw_event}
\end{figure}

The ``Raw Event Data'' packet described in
Figure~\ref{fig:protocol_packet_raw_event} is generated by the neutron
detector system and is used to convey time-of-flight and physical location
information for detected neutrons. It contains the following fields:
\begin{itemize}
\item{\bf Length} is defined by the number of events in the packet. There
are 16 bytes of pulse and packet information followed by 8 bytes of data
for each event.
\item{\bf Pulse ID} identifies the most recent pulse generated by the
accelerator. This ID may not be orginating pulse for the events in this packet
due to frame skew.
\item{\bf Source ID} identifies the upstream detector hardware that
processed these events.
\item{\bf Cycle} identifies the Accelerator cycle number for the most
resent pulse. This may not be the originating pulse for the events in this
packet due to frame skew.
\item{\bf Pkt Seq} is a monotonically increasing event packet number for the
current pulse.
\item{\bf EOP} (End of Packet) indicates that this is the last event packet
to be sent for the current pulse.
\item{\bf Pulse Charge} indicates the charge of the most recent pulse. This
number must be mulitplied by 10 to get the charge in picoColumbs. This
value may not be the charge for the events in this packet due to frame skew.
\item{\bf Pulse Energy} indicates the voltage of the most recent pulse,
given in electron volts (See TODO). This value may not be the voltage for the
events in this packet due to frame skew.
\item{\bf Event(s)} give the time of flight and location of the detected
neutron(s), in a format given by Figure~\ref{fig:protocol_packet_event_data}.
Pixel ID is the physical pixel associated with this neutron event.
\end{itemize}

When run over a network transport that does not provide reliable service, such
as UDP, the protocol uses the {\bf Pulse ID}, {\bf Source ID}, {\bf Pkt Seq},
and {\bf EOP} field to detect the loss of packets describing the events for a
pulse. Each packet source ({\bf Source ID}) sequentially numbers the raw event
packets it sends for a pulse, starting at zero. The source must set {\bf EOP}
on the last raw event packet it sends for a pulse before sending any packet
with a new {\bf Pulse ID}. The receiver will then be able to detect lost
packets by the missing sequence numbers. Should the {\bf EOP} packet be lost,
this will be detected when the sources sends the first packet with a new {\bf
Pulse ID}. This loss-detection protocol assumes that the transport does not
reorder packets; additional measures must be taken in the receiver should be
taken to deal with reordering should it be necessary.

There are seven unused bits ({\bf Reserved}) for this packet type. These bits
must be zero.

``Raw Event Data'' packets may also be used to transmit ``fast'' metadata
about the sample environment, such as the torque being applied to the sample
by a load frame. These packets will be distinguishable from neutron data
by theour {\bf Source ID} and may also feature unique pixel IDs that further
differentiate them from neutron event packets.

\todo{Define units for ``Pulse Charge''}{%
There is inconsistent documentation about the units for ``Pulse Charge''. It
either in units of 10 nC or pC. Steve Hartman is checking this.}
\todo{Define units for ``Pulse Energy''}{%
For the raw event data, the RTDL provides a ``ring period'' value which is
the time it takes the proton beam to travel around the fixed circumference
of the accumulator ring. The field is in picoseconds (Check with Steve
Hartman to verify.) Should the raw event data just copy that value to here
and let downstream deal with the translation?}
\todo{Better documentaion of frame skew}{%
It is discussed in ``Neutron Event Timing'', but we need to be specific
about it anywhere it could be confused.}

\begin{figure}
  \centering
  \begin{bytefield}{32}
    \bitheader{31,0} \\
    \wordbox{1}{Time-of-flight (units of 100ns, u32)} \\
    \wordbox{1}{Pixel ID (u32)}
  \end{bytefield}
  \caption{Event Data Format}
  \label{fig:protocol_packet_event_data}
\end{figure}

\newpage
\subsubsection{Detector System: Real Time Data Link Information}

\begin{figure}[h]
  \centering
  \begin{bytefield}{32}
    \bitheader{31,0} \\
    \wordbox{1}{64} \\
    \wordbox{1}{TBD} \\
    \wordbox{1}{Pulse ID (seconds)} \\
    \wordbox{1}{Pulse ID (nanoseconds)} \\

    \bitheader{31,0} \\
    \wordbox{1}{Pulse Charge (units of 10 pC, u32)} \\
    \wordbox{1}{Ring Period (picoseconds, u32)} \\
    \wordbox[lrt]{1}{14 Reserved words} \\
    \skippedwords \\
    \wordbox[lrb]{1}{}
  \end{bytefield}
  \caption{Packet Type TBD: Real Time Data Link Information}
  \label{fig:protocol_packet_rtdl}
\end{figure}

The ``Real Time Data Link'' packet described in
Figure~\ref{fig:protocol_packet_rtdl} is generated by the neutron
detector system and is used to information about the most recent pulse
from the accelerator.  It contains the following fields:
\begin{itemize}
\item{\bf Length} is always 64 bytes.
\item{\bf Pulse ID} identifies the most recent pulse generated by the
accelerator.
\item{\bf Pulse Charge} indicates the charge of the most recent pulse units of
10 picoColumbs. This value may not be the charge for the events in this packet
due to frame skew.
\item{\bf Ring Period} indicates the number of picoseconds the proton beam
takes to travel the circumference of the accumulator ring. This is used
to calculate beam energy.
\item{\bf Reserved} is 52 bytes of room for future values from the RTDL
system (see TODO).
\end{itemize}

\todo{Accelerator Cycle in RTDL?}{%
We have the accelerator cycle number in the raw event packet, but not in
the RTDL packet?}
\todo{RTDL Reserved fields?}{%
Do we need to have reserved fields in the RTDL packet? Once a packet format
has been seen in production, we're not supposed to change its definition,
so how would follow on software know if these fields are now valid vs
having a zero in them when processing historical event files?}

\newpage
\subsection{Stream Management Service generated events}

The following packet types are generally generated by, or used when
communicating with the Stream Management Service. These packets generally use
the current time-of-day clock for their {\bf Timestamp} fields, though the
Banked Event Data packets will use the Pulse ID. While all of these
packets may be seen over a network transport, several have meaning only
when seen on permament storage and should be ignored when seen on the network.


\newpage
\subsubsection{Banked Event Data}

\begin{figure}[h]
  \centering
  \begin{bytefield}{32}
    \bitheader{31,0} \\
    \wordbox{1}{Length $8 N + 16$} \\
    \wordbox{1}{TBD} \\
    \wordbox{1}{Pulse ID (seconds)} \\
    \wordbox{1}{Pulse ID (nanoseconds)} \\

    \bitheader{31,16,15,0} \\
    \wordbox{1}{Pulse Charge (units of 10 pC, u32)} \\
    \wordbox{1}{Pulse Energy (eV, u32)} \\
    \wordbox{1}{Ring Period (picoseconds, u32)} \\
    \wordbox{1}{Accelerator Cycle} \\
    \bitbox{16}{Flags (u16)} &
    \bitbox{16}{Bank ID (u16)} \\
    \wordbox{2}{Event 1} \\
    \wordbox{1}{$\cdots$} \\
    \wordbox{2}{Event $N$}
  \end{bytefield}
  \caption{Packet Type TBD: Banked Event Data}
  \label{fig:protocol_packet_banked_event_data}
\end{figure}

The ``Banked Event Data'' packet described in
Figure~\ref{fig:protocol_packet_banked_event_data} is generated by the
Stream Management Service and is used to convey time-of-flight and logical
location information for detected neutrons. It contains the following fields:
\begin{itemize}
\item{\bf Length} is defined by the number of events in the packet. There
are 16 bytes of pulse and packet information followed by 8 bytes of data
for each event.
\item{\bf Pulse ID} identifies the most recent pulse generated by the
accelerator. This ID may not be orginating pulse for the events in this packet
due to frame skew.
\item{\bf Cycle} identifies the Accelerator cycle number for the most
resent pulse. If this value is unknown (due to an error such as a lost RTDL
packet), this field will be set to all ones. This may not be the originating
pulse for the events in this packet due to frame skew.
\item{\bf Pulse Charge} indicates the charge of the most recent pulse in units
of 10 picoColumbs. This value may not be the charge for the events in this
packet due to frame skew.
\item{\bf Pulse Energy} indicates the voltage of the most recent pulse,
given in electron volts. This value is calculated from the {\bf Ring Period}
field and depends on the accumulator ring geometry. This value may not be the
voltage for the events in this packet due to frame skew.
\item{\bf Ring Period} indicates the number of picoseconds the proton beam
takes to travel the circumference of the accumulator ring. This is used
to calculate beam energy.
\item{\bf Bank ID} indicates which detector bank is the source of the events
in this packet.
\item{\bf Flags} gives space for indicating various properties about these
events (see TODO).
\item{\bf Event(s)} give the time of flight and location of the detected
neutron(s), in a format given by Figure~\ref{fig:protocol_packet_event_data}.
Pixel ID is the logical pixel, translated from the physical ID given by the
detector.
\end{itemize}

While the ``Raw Event Data'' packet described in
section~\ref{section:protocol_raw_event_packet} is defined for primary use over
an unreliable transports and relatively small transmission sizes, the ``Banked
Event Data'' packet is oriented for consumption by analysis systems
communicating over reliable transports, with the ability to send larger
messages. Thus, the Banked packet does not require mechanisms to detect lost
events for a pulse.

The physical layout of pixel IDs in the detector systems for an instrument
do not always match up with the layout that the scientists wish to use for
analysis. The Stream Managment Service will transform the pixel IDs from
physical in the Raw packets to logical in the Banked packets using a
one-to-one mapping (communicated in the stream via ``Pixel Mapping Table''
packets defined in Section~\ref{section:protocol_pixel_mapping_table}.
These logical pixels are then communicated in the ``Banked Event Data''
packets. As part of this mapping process, the detector bank is determined,
and the events are sorted into buckets accordingly.

\todo{Define Banked Event Data Flags field}{%
We have space here, are there flags that make sense in this context?}

\newpage
\subsubsection{Beam Monitor Event Data}

\begin{figure}[h]
  \centering
  \begin{bytefield}{32}
    \bitheader{31,0} \\
    \wordbox{1}{Length $4 N + 16$} \\
    \wordbox{1}{TBD} \\
    \wordbox{1}{Pulse ID (seconds)} \\
    \wordbox{1}{Pulse ID (nanoseconds)} \\

    \bitheader{31,16,15,0} \\
    \wordbox{1}{Pulse Charge (units of 10 pC, u32)} \\
    \wordbox{1}{Pulse Energy (eV, u32)} \\
    \wordbox{1}{Ring Period (picoseconds, u32)} \\
    \wordbox{1}{Accelerator Cycle} \\
    \bitbox{16}{Flags (u16)} &
    \bitbox{16}{Monitor ID (u16)} \\
    \wordbox{1}{Time-of-flight 1 (units of 100 ns, u32)} \\
    \wordbox{1}{$\cdots$} \\
    \wordbox{1}{Time-of-flight $N$} \\
  \end{bytefield}
  \caption{Packet Type TBD: Beam Monitor Event Data}
  \label{fig:protocol_packet_monitor_event_data}
\end{figure}

The ``Beam Monitor Event Data'' packet described in
Figure~\ref{fig:protocol_packet_monitor_event_data} is generated by the
Stream Management Service and is used to convey neutron time-of-flight
information originating from beam monitors. It contains the following fields:
\begin{itemize}
\item{\bf Length} is defined by the number of events in the packet. There
are 16 bytes of pulse and packet information followed by 4 bytes of data
for each event.
\item{\bf Pulse ID} identifies the most recent pulse generated by the
accelerator. This ID may not be orginating pulse for the events in this packet
due to frame skew.
\item{\bf Cycle} identifies the Accelerator cycle number for the most
resent pulse. If this value is unknown (due to an error such as a lost RTDL
packet), this field will be set to all ones. This may not be the originating
pulse for the events in this packet due to frame skew.
\item{\bf Pulse Charge} indicates the charge of the most recent pulse in units
of 10 picoColumbs. This value may not be the charge for the events in this
packet due to frame skew.
\item{\bf Pulse Energy} indicates the voltage of the most recent pulse,
given in electron volts. This value is calculated from the {\bf Ring Period}
field and depends on the accumulator ring geometry. This value may not be the
voltage for the events in this packet due to frame skew.
\item{\bf Ring Period} indicates the number of picoseconds the proton beam
takes to travel the circumference of the accumulator ring. This is used
to calculate beam energy.
\item{\bf Monitor ID} indicates which beam monitor is the source of the events
in this packet.
\item{\bf Flags} gives space for indicating various properties about these
events (see TODO).
\item{\bf Time-of-flight} give the time-of-flight in units of 100 nanoseconds
of the detected neutrons.
\end{itemize}

While the ``Raw Event Data'' packet described in
section~\ref{section:protocol_raw_event_packet} is defined for primary use over
an unreliable transports and relatively small transmission sizes, the ``Beam
Monitor Event Data'' packet is oriented for consumption by analysis systems
communicating over reliable transports, with the ability to send larger
messages. Thus, the Beam Monitor packet does not require mechanisms to detect
lost events for a pulse.

Beam monitors do not record pixel IDs to record the location of detected
neutrons, so there is no need to spend the space storing that information.
This is especially important as the monitors may often be exposed to ``white
beam'' conditions, generating non-trivial amounts of data in spite of their
low detection efficiency.

\todo{Define Banked Event Data Flags field}{%
We have space here, are there flags that make sense in this context?}

\newpage
\subsubsection{Pixel Mapping Table Packet}
\label{section:protocol_pixel_mapping_table}

\begin{figure}[h]
  \centering
  \begin{bytefield}{32}
    \bitheader{31,0} \\
    \wordbox{1}{$N$} \\
    \wordbox{1}{TBD} \\
    \wordbox{1}{Timestamp (seconds)} \\
    \wordbox{1}{Timestamp (nanoseconds)} \\

    \bitheader{31,0} \\
    \wordbox{2}{Mapping Data ($N$ bytes)} \\
  \end{bytefield}
  \caption{Packet Type TBD: Pixel Map Packet}
  \label{fig:protocol_packet_pixel_map}
\end{figure}

As part of its processing of the neutron event data, the SMS will map the
physical pixel ID to a logical pixel ID. This mapping must be one-to-one and
reverisble. The ``Pixel Mapping Table'' packet defined in
Figure~\ref{fig:protocol_packet_pixel_map} will preserve the table for
consumers of the event stream. This packet must be present near the beginning
of a new experiment recording, but may occur multiple times in the transmitted
stream. Each additional copy transmitted must be identical to the first copy
for the same experiment recording.

\begin{figure}[ht]
  \centering
  \begin{bytefield}{48}
    \bitheader{48,16,15,0} \\
    \bitbox{32}{Base Physical ID (u32)} &
    \bitbox{16}{Count (u16)} \\
    \wordbox{1}{Map Entry 0} \\
    \wordbox[lrt]{1}{Map Entry 1} \\
    \wordbox[lr]{1}{$\cdots$} \\
    \wordbox[lrb]{1}{Map Entry $N$} \\
  \end{bytefield}
  \caption{Pixel Map Section}
  \label{fig:protocol_packet_pixel_map_section}
\end{figure}

\begin{figure}[ht]
  \centering
  \begin{bytefield}{48}
    \bitheader{48,16,15,0} \\
    \bitbox{32}{Logical ID (u32)} &
    \bitbox{16}{Bank ID (u16)} \\
  \end{bytefield}
  \caption{Pixel Map Entry}
  \label{fig:protocol_packet_pixel_map_entry}
\end{figure}

The physical pixel ID spase is not always compactly populated. It often
contains sequential clusters of IDs with large gaps between the clusters.
To reduce the size of the mapping table, it is described using sections as
defined in Figure~\ref{fig:protocol_packet_pixel_map_section}.
{\bf Base Physical ID} defines be offset into the physical ID space for
which this section defines mappings. {\bf Count} gives the number of
mapping entries in this section.

The mapping entries are defined in
Figure~\ref{fig:protocol_packet_pixel_map_entry}.  The physical pixel ID for
the entry is determined by adding its zero-based index into the section to the
{\bf Base Physical ID}. This physical ID is then mapped to the {\bf Logical ID}
specified, and was generated by a detector in bank {\bf Bank ID}.



\newpage
\subsubsection{Synchronization Packet}

\begin{figure}[h]
  \centering
  \begin{bytefield}{32}
    \bitheader{31,0} \\
    \wordbox{1}{28 + roundup($N$, 4)} \\
    \wordbox{1}{TBD} \\
    \wordbox{1}{Timestamp (seconds)} \\
    \wordbox{1}{Timestamp (nanoseconds)} \\

    \bitheader{31,0} \\
    \wordbox{4}{Signature (16 bytes)} \\
    \wordbox{2}{File Offset (u64)} \\
    \wordbox{1}{Comment Length (u32)} \\
    \wordbox{2}{Comment ($N$ byte string, optional)} \\
  \end{bytefield}
  \caption{Packet Type TBD: Synchronization Packet}
  \label{fig:protocol_packet_sync}
\end{figure}

While the ADARA protocol is primarily used for network transmission of
event data, it is necessary to store the raw data on storage media during
processing and for archival storage. The ``Synchronization`` packet defined
in Figure~\ref{fig:protocol_packet_sync} allows partial recovery of experiment
data in the event that a media error or other corruption destroys a packet
header. Without a marker in the data stream, it may be impossible to find
the beginning of the next packet, rendering all data after the media defect
unusable. Periodically placing synchronization packets in the stream contains
the damage that the error can cause.

\begin{itemize}
\item{\bf Length} is 28 bytes plus the length of the optional comment field,
rounded up to the nearest multiple of 4 bytes.
\item{\bf Timestamp} is the current system time when the packet is stored
to disk.
\item{\bf Signature} is the byte string 0x53 0x4e 0x53 0x41 0x44 0x41 0x52
0x41 0x4f 0x52 0x4e 0x4c 0x00 0x00 0xf0 0x7f.
\item{\bf File Offset} is the offset from the beginning of the file to the
length field of this packet. This is used to further verify that we have a
valid Synchronization Packet.
\item{\bf Comment Length} is the length of the comment that follows, not
including padding.
\item{\bf Comment} is an optional UTF-8 encoded comment. This may be used
to store a small amount of human readable information that may be displayed
by the Unix `file' command. This field is padded by zeros to the nearest
multiple of 4.
\end{itemize}

Every file containing ADARA data should begin with a Synchronization packet to
allow system utilities to identify the file as such. Additional synchronization
packets should be added to the file as desired to contain any potential damage.
As the Synchonization packet is only useful when used for storage, it should be
ignored when encountered over a network transport.

{\bf Signature} was constructed using the string ``SNSADARA'' and concatenating
a positive, signaling NaN value in IEEE 754 double format. The mantessa bits
marked as ``don't care'' in the standard have the string ``ORNL'' inserted such
that it views properly when displayed in byte order.  As it is expected that
NaN values are unlikely to be part of a valid data stream, this signature
should prove unique when combined with a matching file offset.


\newpage
\subsubsection{Run Status Packet}

\begin{figure}[h]
  \centering
  \begin{bytefield}{32}
    \bitheader{31,0} \\
    \wordbox{1}{12} \\
    \wordbox{1}{TBD} \\
    \wordbox{1}{Timestamp (seconds)} \\
    \wordbox{1}{Timestamp (nanoseconds)} \\

    \bitheader{31,24,23,0} \\
    \wordbox{1}{Run Number(u32)} \\
    \wordbox{1}{Run Start (seconds, u32)} \\
    \bitbox{8}{Status} &
    \bitbox{24}{File Number} \\
  \end{bytefield}
  \caption{Packet Type TBD: Run Status Packet}
  \label{fig:protocol_packet_status}
\end{figure}

\begin{table}[h]
  \begin{center}
    \begin{tabular}{c | l}
	Value & Definition \\
	\hline
	0 & No current run \\
	1 & Beginning new run \\
	2 & Continuing current run (new file marker) \\
    \end{tabular}
  \end{center}
  \caption {Run Status Definitions}
  \label{table:protocol_run_status_values}
\end{table}

The ``Run Status'' packet is used to indicate if we are currently in an
experimental run, and if so, identify it. This packet must be sent/written at
the start of each experimental run, and should be written to the start of
each event data file to tag it with the run data it contains.

\begin{itemize}
\item{\bf Length} is 12 bytes.
\item{\bf Timestamp} is the current system time when the packet is created.
\item{\bf Run Number} is the identifier for the current experimental run. This
field is 0 if we are not currently recording an experimental run.
\item{\bf Run Start} contains the seconds portion of the time that we began
the current experiment run. This field is 0 if we are not currently recording
data.
\item{\bf Status} is defined in Table~\ref{table:protocol_run_status_values}.
\item{\bf File Number} indentifies which file this data is stored in. This
field is undefined in we are not currently recording an experimental run.
\end{itemize}

The Stream Management Service breaks experimental recordings into multiple
chunks to allow a gradual reclaimation of storage space once the recording
has been translated to a NeXus file. This process increases the window in
which at least partial data from a experiment can be retrieved from the SMS
in the case that a problem is detected after the guaranteed recovery period.
Rather than relying solely on the filename (which could be subject to human
error) to determine how much of the recording has been preserved, the
``Run Status'' packet provides an authoratative source for the amount
of time and approximate data missing by using the {\bf Run Start} and
{\bf File Number} fields along with the administrator controlled chunk size.

\newpage
\subsubsection{Heartbeat Packet}

\begin{figure}[h]
  \centering
  \begin{bytefield}{32}
    \bitheader{31,0} \\
    \wordbox{1}{0} \\
    \wordbox{1}{TBD} \\
    \wordbox{1}{Timestamp (seconds)} \\
    \wordbox{1}{Timestamp (nanoseconds)}
  \end{bytefield}
  \caption{Packet Type TBD: Heartbeat Packet}
  \label{fig:protocol_packet_heartbeat}
\end{figure}

In order for clients to distinguish between a crashed server and one with
no event data to transmit (ie, shutter is closed), there must be periodic
traffic over the network transport. The ``Heartbeat'' packet fills this need.

This role may also be performed by the network transport's keepalive settings,
though that does not help communication with the detector system which is
expected to use an unreliable transport.

\newpage
\subsubsection{Run Information Packet}

\begin{figure}[h]
  \centering
  \begin{bytefield}{32}
    \bitheader{31,0} \\
    \wordbox{1}{4 + roundup($N$, 4)} \\
    \wordbox{1}{TBD} \\
    \wordbox{1}{Timestamp (seconds)} \\
    \wordbox{1}{Timestamp (nanoseconds)} \\

    \bitheader{31,0} \\
    \wordbox{1}{XML Length (u32)} \\
    \wordbox{4}{XML Run Information ($N$ byte string data)} \\
  \end{bytefield}
  \caption{Packet Type TBD: Run Information Packet}
  \label{fig:protocol_packet_runinfo_run}
\end{figure}

\todo{Define Run Information packet}{%
We need to finalize the mechanims for getting this information to interested
parties. Jim and Marie will define the XML schema for this packet, and we
will reference it from the protocol document. User comments, bin sizes,
and other information that is not specifically related will be moved to
metadata value packets.}


\newpage
\subsubsection{Client Hello Packet}

\begin{figure}[h]
  \centering
  \begin{bytefield}{32}
    \bitheader{31,0} \\
    \wordbox{1}{4} \\
    \wordbox{1}{TBD} \\
    \wordbox{1}{Timestamp (seconds)} \\
    \wordbox{1}{Timestamp (nanoseconds)} \\

    \bitheader{31,0} \\
    \wordbox{1}{Requested Start Timestamp (seconds, u32)} \\
  \end{bytefield}
  \caption{Packet Type TBD: Client Hello Packet}
  \label{fig:protocol_client_hello}
\end{figure}

After first connecting to the Stream Management Service, a client will send
a ``Client Hello'' packet to indicate from where in the stream it would like
begin receiving event data. No data will be sent to the client prior to
receiving this packet.

\begin{itemize}
\item{\bf Length} is 4 bytes.
\item{\bf Timestamp} is the current system time when the packet is created.
\item{\bf Requested Start Timestamp} suggests to the server where in the live
stream to start sending events to this client, as defined below.
\end{itemize}

The value given in {\bf Requested Start Timestamp} is a suggestion to the
server and may be ignored. The client may receive events preceeding its
requested timestamp, or may only receive live data if historical data is not
available.

The requested value is bounded by run state transitions. If the client
requests data from before the current run started, the server will start
sending events from the beginning of the run. If there is no current run and
the requested timestamp is before the end of the previous run, the server will
begin sending events from immediately after the previous run ended.

The requested timestamp may be rounded to the granularity of the event indexing
used by the Stream Management Service. In this case, the timestamp will be
rounded towards the historical data, and the client may receive events that
preceed the requested point in time. The client is responsible for filtering
the undesired requests from the the stream.

If the {\bf Requested Start Timestamp} is zero, the server will not send
historical data and will begin sending events with the next pulse.

\newpage
\subsubsection{Run Completion Packet}

\begin{figure}[h]
  \centering
  \begin{bytefield}{32}
    \bitheader{31,0} \\
    \wordbox{1}{4 + roundup($N$, 4)} \\
    \wordbox{1}{TBD} \\
    \wordbox{1}{Timestamp (seconds)} \\
    \wordbox{1}{Timestamp (nanoseconds)} \\

    \bitheader{31,16,15,0} \\
    \bitbox{16}{Status (u16)} &
    \bitbox{16}{Reason Length (u16)} \\
    \wordbox{2}{Reason ($N$ byte string, optional)} \\
  \end{bytefield}
  \caption{Packet Type TBD: Run Completion Packet}
  \label{fig:protocol_packet_run_completion}
\end{figure}

\begin{table}[h]
  \begin{center}
    \begin{tabular}{c | l}
	Value & Definition \\
	\hline
	0 & Success \\
	0x0001 to 0x7fff & Transient Error \\
	0x8000 to 0xffff & Permament Error \\
    \end{tabular}
  \end{center}
  \caption {Status Definitions}
  \label{table:protocol_completion_status_values}
\end{table}

The ``Completion'' Packet allows the Stream Management Service to signal to the
Translation Service that it has transmitted all of the data related to an
experimental run. The Translation Servce will use this same packet format to
indicate the results of translation back to the SMS.

\begin{itemize}
\item{\bf Length} is 8 bytes plus the length of the optional {\bf Reason} field
rounded up to the next multiple of four.
\item{\bf Timestamp} is the current system time when the packet is created.
\item{\bf Status} indicates the result of translation as defined in
Table~\ref{table:protocol_completion_status_values}.
\item{\bf Reason Length} is the length of the comment that follows, not
including padding.
\item{\bf Reason} is an optional UTF-8 encoded reason. This may be used
to communicate a small amount of human readable information to explain
a non-successful status. This field is padded by zeros to the nearest
multiple of 4.
\end{itemize}

If the {\bf Status} contains a value defined to be a Transient Error, the SMS
should attempt to resend this experimental run to the Translation Service.  If
the {\bf Status} contains a value defined to be a Permament Error, the SMS
should not resend this experimental run until a human has inspected and
corrected the root cause. Only if {\bf Status} is successful may the SMS
consider the experimental run safely on disk at the Translation Service and
mark the storage space valid for eventual reclaimation.

\newpage
\subsubsection{Statistics Reset Packet}

\begin{figure}[h]
  \centering
  \begin{bytefield}{32}
    \bitheader{31,0} \\
    \wordbox{1}{0} \\
    \wordbox{1}{TBD} \\
    \wordbox{1}{Timestamp (seconds)} \\
    \wordbox{1}{Timestamp (nanoseconds)}
  \end{bytefield}
  \caption{Packet Type TBD: Statistics Reset Packet}
  \label{fig:protocol_packet_stats_reset}
\end{figure}

During manual operation of a beamline such as performing sample alignment, the
scientist/user may not be recording neutron event data for translation. In
these situations, the user may wish to reset the live histogram generation and
any other temporary or partial analysis result between each phase of the
process. To allow ease of use and manipulation, the live tools should be able
to receive this request via the same data stream as the experiment data. This
allows precise, common boundaries between data sets to be aligned among all of
the live tools, and avoids the need to have the user manipulate multiple tools
to perform the reset.

When non-translation clients receive a ``Statistics Reset Packet'' (defined in
Figure~\ref{fig:protocol_packet_stats_reset}) in the data stream, they should
discard previous neutron data. The Translation Service may note the occurance
of these packets in the appropriate place.

The contents of the packet field are as follows:
\begin{itemize}
\item{\bf Length} is 0 bytes.
\item{\bf Timestamp} is the current system time when the packet is created.
\end{itemize}


\newpage
\subsection{Meta Data Packets}

In addition to the neutron events from the detector systems, the SNS must
collect information about the environment of the sample being studied. This
data comes from a diverse and dynamic collection of devices.

\todo{Expand metadata description}{%
Jim has a good description about the metadata, just need to clean it up
and bring it into this document.}
\todo{Define Variable Status/Severity Fields}{%
We should have a common table here for the status and severity fields. This
is likely to be derived from EPICS.}
\todo{Device ID 0 is a pseudo device}{%
We will use device ID 0 as a pseudo device to collect user comments and
parameters such as the binning configuration for the translation service.
There will also be a system parameter to collect ``system status'' information
that is not already handled by variables specific to the sample environment
devices.}
\todo{Fast metadata handling}{%
I'm expecting to translate fast metadata into these packets; we need to
verify that chopper data will fit, and determine exactly how we should
handle it.}

\newpage
\subsubsection{Device Descriptor Packet}

\begin{figure}[h]
  \centering
  \begin{bytefield}{32}
    \bitheader{31,0} \\
    \wordbox{1}{8 + roundup($N$, 4)} \\
    \wordbox{1}{TBD} \\
    \wordbox{1}{Timestamp (seconds)} \\
    \wordbox{1}{Timestamp (nanoseconds)} \\

    \bitheader{31,16,15,0} \\
    \wordbox{1}{Device ID (u32)} \\
    \wordbox{1}{Descriptor Length (u32)} \\
    \wordbox[lrt]{1}{XML Descriptor ($N$ byte string)} \\
    \skippedwords \\
    \wordbox[lrb]{1}{}
  \end{bytefield}
  \caption{Packet Type TBD: Device Descriptor Packet}
  \label{fig:protocol_packet_device_desc}
\end{figure}

\begin{itemize}
\item{\bf Length} is 8 bytes plus the length of the {\bf XML Descriptor} field
rounded up to the next multiple of four.
\item{\bf Timestamp} is the current system time when the packet is created.
\item{\bf Device ID} indicates the device this descriptor describes.
\item{\bf Descriptor Length} is the length of the XML descriptor that follows,
not including padding.
\item{\bf XML Descriptor} is a UTF-8 encoded XML document describing the
properties of this variable. This field is padded by zeros to the nearest
multiple of 4.
\end{itemize}

The value in the {\bf Device ID} field is used to identify a device in the
ADARA, and to provide separate variable namespaces for each one.

\todo{Define Device Descriptor XML Schema}{%
Jim and Marie will define the schema for the XML content of this packet.}

\newpage
\subsubsection{Variable Value (u32) Packet}

\begin{figure}[h]
  \centering
  \begin{bytefield}{32}
    \bitheader{31,0} \\
    \wordbox{1}{12} \\
    \wordbox{1}{TBD} \\
    \wordbox{1}{Timestamp (seconds)} \\
    \wordbox{1}{Timestamp (nanoseconds)} \\

    \bitheader{31,16,15,0} \\
    \wordbox{1}{Device ID (u32)} \\
    \wordbox{1}{Variable ID (u32)} \\
    \bitbox{16}{Status} &
    \bitbox{16}{Severity} \\
    \wordbox{1}{Variable Value (u32)}
  \end{bytefield}
  \caption{Packet Type TBD: Variable Value (u32) Packet}
  \label{fig:protocol_packet_value_u32}
\end{figure}

\begin{itemize}
\item{\bf Length} is 12 bytes.
\item{\bf Timestamp} indicates the time that the variable changed.
\item{\bf Device ID} indicates the device this variable is from.
\item{\bf Variable ID} indicates the variable that changed.
\item{\bf Status} of this variable. See TODO.
\item{\bf Severity} of this variable. See TODO.
\item{\bf Variable Value} is an unsigned 32 bit value for this variable.
\end{itemize}

Note that this packet type is also used to carry enumerated and boolean values.
The specific interpretation of the value is dictated by the Device Descriptor
associated with the Device ID.

\newpage
\subsubsection{Variable Value (double) Packet}

\begin{figure}[h]
  \centering
  \begin{bytefield}{32}
    \bitheader{31,0} \\
    \wordbox{1}{16} \\
    \wordbox{1}{TBD} \\
    \wordbox{1}{Timestamp (seconds)} \\
    \wordbox{1}{Timestamp (nanoseconds)} \\

    \bitheader{31,16,15,0} \\
    \wordbox{1}{Device ID (u32)} \\
    \wordbox{1}{Variable ID (u32)} \\
    \bitbox{16}{Status} &
    \bitbox{16}{Severity} \\
    \wordbox{2}{Variable Value (IEEE 754 double)}
  \end{bytefield}
  \caption{Packet Type TBD: Variable Value (double) Packet}
  \label{fig:protocol_packet_value_double}
\end{figure}

\begin{itemize}
\item{\bf Length} is 16 bytes.
\item{\bf Timestamp} indicates the time that the variable changed.
\item{\bf Device ID} indicates the device this variable is from.
\item{\bf Variable ID} indicates the variable that changed.
\item{\bf Status} of this variable. See TODO.
\item{\bf Severity} of this variable. See TODO.
\item{\bf Variable Value} is an IEEE 754 double floating point value for this
variable.
\end{itemize}

\newpage
\subsubsection{Variable Value (string) Packet}

\begin{figure}[h]
  \centering
  \begin{bytefield}{32}
    \bitheader{31,0} \\
    \wordbox{1}{12 + roundup($N$, 4)} \\
    \wordbox{1}{TBD} \\
    \wordbox{1}{Timestamp (seconds)} \\
    \wordbox{1}{Timestamp (nanoseconds)} \\

    \bitheader{31,16,15,0} \\
    \wordbox{1}{Device ID (u32)} \\
    \wordbox{1}{Variable ID (u32)} \\
    \bitbox{16}{Status} &
    \bitbox{16}{Severity} \\
    \wordbox{1}{String Length (u32)} \\
    \wordbox{2}{Variable Value ($N$ byte string data)} \\
  \end{bytefield}
  \caption{Packet Type TBD: Variable Value (string) Packet}
  \label{fig:protocol_packet_value_string}
\end{figure}

\begin{itemize}
\item{\bf Length} is 12 bytes plus the length of the {\bf Variable Value} field
rounded up to the next multiple of four.
\item{\bf Timestamp} indicates the time that the variable changed.
\item{\bf Device ID} indicates the device this variable is from.
\item{\bf Variable ID} indicates the variable that changed.
\item{\bf Status} of this variable. See TODO.
\item{\bf Severity} of this variable. See TODO.
\item{\bf String Length} is the length of the string that follows,
not including padding.
\item{\bf Variable Value} is a UTF-8 encoded string containing the new value
for this variable. This field is padded by zeros to the nearest multiple of 4.
\end{itemize}


